\documentclass[oneside]{amsart}
\usepackage[all]{xy}
\usepackage{tikz-cd}
\usepackage[T1]{fontenc}
\usepackage{xstring}
\usepackage{xparse}
\usepackage{xr-hyper}
\usepackage{xcolor}
\definecolor{brightmaroon}{rgb}{0.76, 0.13, 0.28}
\usepackage[linktocpage=true,colorlinks=true,hyperindex,citecolor=blue,linkcolor=brightmaroon]{hyperref}
\usepackage[nameinlink]{cleveref}
\usepackage[left=1.25in,right=1.25in,top=0.75in,bottom=0.75in]{geometry}
\setlength{\marginparwidth}{2cm}
%\usepackage[charter,greekfamily=didot]{mathdesign}
%\usepackage{Baskervaldx}
\usepackage{amssymb}
\usepackage{stmaryrd}
\usepackage{mathrsfs}
\usepackage{mathpazo}
\usepackage{eulervm}
\linespread{1.05}

\usepackage[nobottomtitles]{titlesec}
\usepackage{marginnote}
\usepackage{enumerate}
\usepackage{longtable}
\usepackage{aurical}
\usepackage{microtype}
\usepackage{mathtools}

\newtheoremstyle{ega-env-style}%
{}{}{\rmfamily}{}{\bfseries}{.}{ }{\thmnote{(#3)}}%

\newtheoremstyle{ega-thm-env-style}%
{}{}{\itshape}{}{\bfseries}{. --- }{ }{\thmname{#1}\thmnumber{ #2}\thmnote{ (#3)}}%

\newtheoremstyle{ega-defn-env-style}%
{}{}{\rmfamily}{}{\bfseries}{. --- }{ }{\thmname{#1}\thmnumber{ #2}\thmnote{ (#3)}}%

\theoremstyle{ega-env-style}
\newtheorem*{env}{---}

\theoremstyle{ega-thm-env-style}
\newtheorem{theorem}{Theorem}[section]
\newtheorem{proposition}[theorem]{Proposition}
\newtheorem{lemma}[theorem]{Lemma}
\newtheorem{corollary}[theorem]{Corollary}
\newtheorem{conjecture}[theorem]{Conjecture}

\theoremstyle{ega-defn-env-style}
\newtheorem{definition}[theorem]{Definition}
\newtheorem{example}[theorem]{Example}
\newtheorem*{examples}{Examples}
\newtheorem*{remark}{Remark}
\newtheorem*{remarks}{Remarks}
\newtheorem*{notation}{Notation}
\newtheorem*{exercise}{Exercise}
\newtheorem*{properties}{Properties}
\newtheorem*{consequences}{Consequences}
\newtheorem*{note}{Note}
\newtheorem*{fact}{Fact}
\newtheorem*{claim}{Claim}

\makeatletter
\def\l@subsection{\@tocline{2}{0pt}{2.5pc}{2.2pc}{}}
\def
\section{\@startsection{section}{1}%
	\z@{.7\linespacing\@plus\linespacing}{.5\linespacing}%
{\normalfont\bfseries\Large\scshape\centering}}
\renewcommand{\@seccntformat}[1]{%
	\ifnum\pdfstrcmp{#1}{section}=0\textsection\fi%
\csname the#1\endcsname.~}
\makeatother

\def\mathcal{\mathscr}
%% Fonts

\def\sh{\mathcal}                   % sheaf font
\def\bb{\mathbf}                    % bold font
\def\cat{\mathsf}                   % category font
\def\numfont{\mathbb}

%% Font Letters

\def\CA{\mathcal{A}}
\def\CE{\mathcal{E}}
\def\CF{\mathcal{F}}
\def\CG{\mathcal{G}}
\def\CL{\mathcal{L}}
\def\OO{\mathcal{O}}
\def\CX{\mathcal{X}}
\def\b1{\mathbb{1}}

%% Cohomology

\def\CH{\mathrm{H}}                 % cohomology H
\def\CHH{\check{\HH}}               % Čech cohomology H
\def\RD{\mathrm{R}}                 % right derived R
\def\LD{\mathrm{L}}                 % left derived L
\def\dual#1{{#1}^\vee}              % dual
\def\Tor{\operatorname{Tor}}        % Tor
\def\Ext{\operatorname{Ext}}        % Ext
\def\HdR{\mathrm{H}_{\mathrm{dR}}}  % de Rham cohomology
\def\Zc{\underline{\mathbb{Z}}}     % constant sheaf with integer coeffs

%% Categories

\def\A{\cat{A}}                     % category A, usually abelian
\def\C{\cat{C}}                     % category C
\def\op{^\cat{op}}                  % opposite category
\def\Set{\cat{Sets}}                % category of sets
\def\Grp{\cat{Gps}}                 % category of groups
\def\Alg{\cat{Alg}}                 % category of algebras
\def\QCoh{\cat{QCoh}}               % category of quasicoherent sheaves
\def\supertilde{{\,\widetilde{\,}\,}}   % use \supertilde instead of ^\sim
\def\Map{\operatorname{Map}}        % morphisms (any category)
\def\Hom{\operatorname{Hom}}        % morphisms (additive category)
\def\iHom{\underline{\Hom}}         % interal hom
\def\End{\operatorname{End}}        % endomorphisms
\def\Aut{\operatorname{Aut}}        % automorphisms
\def\ker{\operatorname{ker}}        % kernel
\def\img{\operatorname{im}}         % image
\def\coker{\operatorname{coker}}    % cokernel
\def\pr{\operatorname{pr}}          % projection
\def\EssIm{\operatorname{EssIm}}    % essential image
\DeclareMathOperator*{\colim}{colim}   % colimit

\def\DAet{\cat{DA}\et}
\def\SH{\cat{SH}}
\def\Sh{\cat{Sh}}                   % category of sheaves
\def\PSh{\cat{PSh}}                 % category of presheaves

%% Schemes

\def\Proj{\operatorname{Proj}}      % Proj
\def\Supp{\operatorname{Supp}}      % support
\def\Spec{\operatorname{Spec}}      % Spec
\def\Spf{\operatorname{Spf}}        % formal Spec
\def\Aff{\mathbb A}                 % affine space
\def\P{\mathbb{P}}                  % projective space
\def\Pic{\operatorname{Pic}}        % Picard group

%% Standard Operators

\def\codim{\operatorname{codim}}    % codimension
\def\id{\operatorname{id}}          % identity

%% Arrows
\renewcommand{\to}{\mathchoice{\longrightarrow}{\rightarrow}{\rightarrow}{\rightarrow}}
\newcommand{\from}{\mathchoice{\longleftarrow}{\leftarrow}{\leftarrow}{\leftarrow}}
\let\mapstoo\mapsto
\renewcommand{\mapsto}{\mathchoice{\longmapsto}{\mapstoo}{\mapstoo}{\mapstoo}}
\def\isoto{\simeq}                  % isomorphism
\def\simto{\xrightarrow{\sim}}      % isomorphism arrow
\def\surjto{\twoheadrightarrow}     % surjetion
\def\injto{\hookrightarrow}         % injection

%% Under/Over Accents

\newcommand{\wh}[1]{\widehat{#1}}   % hat
\newcommand{\wt}[1]{\widetilde{#1}}    % tilde
\def\ul{\underline}                % underline

%% Spaces

\def\Z{\numfont Z}                  % integers
\def\Q{\numfont Q}                  % rationals
\def\N{\numfont N}                  % naturals
\def\R{\numfont R}                  % reals
\def\H{\mathcal H}

%% Groups

\def\GL{\bb{GL}}                   % general linear group
\def\SL{\bb{SL}}                   % special linear group
\def\Sp{\bb{Sp}}
\def\PSp{\bb{PSp}}
\def\GSp{\bb{GSp}}
\def\det{\operatorname{det}}       % determinant

%% Sub/Superscripts

\def\et{^\text{\'et}}              % \'etale
\def\an{^\text{an}}                % analytic
\def\th{^\text{th}}                % th

%% Misc

\newcommand{\defn}[1]{\textbf{#1}}  % definition highlighting
\def\defeq{:=}                     % definition equation
\def\eps{\varepsilon}              % correct epsilon
\newcommand{\gen}[1]{\left\langle\!\left\langle #1 \right\rangle\!\right\rangle}


\def\shHom{\sh{H}\textup{\kern-2.2pt{\Fontauri\slshape om}}}   % sheaf Hom
\def\shProj{\sh{P}\textup{\kern-2.2pt{\Fontauri\slshape roj}}} % sheaf Proj
\def\shExt{\sh{E}\textup{\kern-2.2pt{\Fontauri\slshape xt}}}   % sheaf Ext
\def\shGr{\sh{G}\textup{\kern-2.2pt{\Fontauri\slshape r}}}     % sheaf Gr
\def\shDer{\sh{D}\,\textup{\kern-2.2pt{\Fontauri\slshape er}}} % sheaf Der
\def\shDiff{\sh{D}\,\textup{\kern-2.2pt{\Fontauri\slshape if{}f}}\,} % sheaf Diff
\def\shHomcont{\sh{H}\textup{\kern-2.2pt{\Fontauri\slshape om.\,cont}}}   % sheaf Hom.cont
\def\shAut{\sh{A}\textup{\kern-2.2pt{\Fontauri\slshape ut}}}   % sheaf Aut
\def\shSym{\sh{S}\textup{\kern-2.2pt{\Fontauri\slshape ym}}}   % sheaf Sym

\makeatletter
\newcommand{\cbigoplus}{\DOTSB\cbigoplus@\slimits@}
\newcommand{\cbigoplus@}{\mathop{\widehat{\bigoplus}}}
\makeatother

\newcommand{\muT}{\mu_{\mathrm{Tam}}}
\newcommand{\F}{\mathbb{F}}
\DeclareMathOperator{\Bun}{Bun}
\DeclareMathOperator{\Tr}{Tr}
\DeclareMathOperator{\Frob}{Frob}
\DeclareMathOperator{\Ran}{Ran}
\newcommand{\red}{_{\mathrm{red}}}
\newcommand{\km}{\mathfrak{m}}
\DeclareMathOperator{\Sym}{Sym}
\DeclareMathOperator{\Shv}{Shv}
\DeclareMathOperator{\Gr}{Gr}
\def\D{\mathbb{D}}
\newcommand{\Sch}{\cat{Sch}}
\newcommand{\Fins}{\cat{Fin}^{\mathrm{s}}}
\newcommand{\Fun}{\cat{Fun}}
\DeclareMathOperator{\ins}{ins}

\mathchardef\mhyphen="2D
\newcommand{\iCat}{\cat{\infty\mhyphen Cat}}
\newcommand{\iCatSt}{\cat{\infty\mhyphen Cat}^{\mathrm{St}}}
\newcommand{\lMod}{\Lambda\cat{\mhyphen Mod}}

\def\surjfrom{\twoheadleftarrow}

\usepackage[backend=biber,style=alphabetic]{biblatex}
\addbibresource{pwg.bib}

\title{The Ran Space}
\author{Micha{\l} Mruga{\l}a}

\begin{document}
\maketitle

\setcounter{section}{0}
\section{Coventions and Notation}
\begin{itemize}
	\item $k$ is an algebraically closed field.
	\item $\Sch$ is the category of separated schemes of finite type over $k$.
	\item A \defn{presheaf} on $\Sch$ is a functor $\Sch\op\to \cat{Set}$. We write $\PSh(\Sch)$ for the category of presheaves of schemes.
	\item $\Fins$ is the category of finite-nonempty sets and surjective maps.
	\item We ignore set theoretic issues.
\end{itemize}

\section{\texorpdfstring{The Ran Space and $\ell$-adic Sheaves}{The Ran Space and l-adic Sheaves}}
\subsection{The Ran Space}
\begin{definition}	
	Let $X$ be a separated scheme. The \defn{Ran space} of $X$ is the presheaf with values
	\[
		(\Ran X)(S) = \left\{ I\subset X(S) \middle| I \in \Fins \right\}.
	\] 
\end{definition}
We will show later that we can view $\Ran X$ as a colimit of schemes:
\[
\Ran X = \colim_{\CI\in (\Fins)\op} X^{\CI}
\] 
Since $X$ is separated, the morphisms $X^{\CJ}\to X^{\CI}$ for a morphism $\CI\to \CJ$ in $\Fins$ are closed embeddings. So $Ran X$ is similar to an ind-scheme, the difference being that the colimit is not filtered.

This description of the Ran space suggests an answer to:

\vspace{0.5em}
\begin{center}
	\textbf{Question:} What is the right notion of an $\ell$-adic sheaf on the Ran space?
\end{center}
\vspace{0.5em}

It should be a colimit of $\ell$-adic sheaves on $X^{\CI}$. But we would like to do more, we would like to make sense of the entire category of sheaves on $\Ran X$ as a colimit of categories of sheaves on $X^{\CI}$.This is where higher category theory comes in.

\subsection{\texorpdfstring{$\infty$-category of $\ell$-adic Sheaves}{Infinity-category of l-adic Sheaves}}
In the last talk, to each scheme we associated the derived category $D(X,\Lambda)$ of $\ell$-adic sheaves with coefficients in
\[
\Lambda\in \left\{ \Z /\ell^{n}, \Z_\ell, \Q_\ell, \overline{\Q}_\ell \middle| \ell\neq \car k \right\} .
\] 
The compact objects are constructible sheaves and they generate $D(X,\Lambda)$ under filtered colimits.

We adapt this definition to the world of $\infty$-categories as follows:
\begin{enumerate}[(i)]
	\item Let $\Shv_c(X)$ be the stable $\infty$-category of constructible $\ell$-adic sheaves with coefficients in $\Lambda$.
	\item Let $\Shv(X)$ be the ind-completion of $\Shv_c(X)$. We call it the \defn{$\infty$-category of $\ell$-adic sheaves} on $X$.
\end{enumerate}

In fact, we have a functor
\begin{align*}
	\Shv:\Sch\op &\longrightarrow \iCat \\
	S &\longmapsto \Shv(S) \\
	(f:S\from S') &\longmapsto (f^{*}:\Shv(S)\to \Shv(S')).
\end{align*}

We want to make sense of the expression
\[
\Shv(\Ran X) \defeq \colim_{\CI\in (\Fins)\op}\Shv(X^{\CI}).
\] 

\subsection{Prestacks}
Instead of just settling on a definition of $\ell$-adic sheaves for $\Ran X$ we will define it for a general (pre)stack. Recall that a stack is a category fibered in groupoids over $\Sch$ satisfying some descent conditions. If we ignore the descent conditions this is essentially the same thing as a pseudofunctor
\[
	\Sch\op \to \cat{Grpds}
\] 
to the 2-category of groupoids. Making this homotopy coherent we arrive at:
\begin{definition}
	A \defn{prestack} to be a contravariant functor
	\[
	\Sch\op\to \Spc
	\] 
	from the category of schemes to the $\infty$-category of spaces (aka $\infty$-groupoids). We write $\PreStk$ for the $\infty$-category of prestacks.
\end{definition}
We have fully faithful functors
\[
\Sch \injto \PSh(\Sch) \injto \PreStk
\] 
the latter preserving limits and filtered colimits. Now we can make sense of:

\begin{proposition}\label{prop:rancolim}
	The Ran space is a colimit of
	\begin{align*}
		F: (\Fins)\op &\longrightarrow \PreStk \\
		\CI &\longmapsto X^{\CI}
	\end{align*}
	via
	\begin{align*}
		\ins_\CI(S): X^{\CI}(S) &\longrightarrow (\Ran X)(S) \\
		f &\longmapsto \img(f).
	\end{align*}
\end{proposition}
\begin{proof}
	We claim that for a set $A$
	\[
	\Ran A = \left\{ I\subset A \middle| I\in \Fins \right\} \isoto \colim_{\CI\in \Fins}A^{\CI}
	\] 
	in the category of spaces. Since colimits in $\PreStk$ are computed pointwise, the result follows by taking $A=X(S)$.

	By \cite[\href{https://kerodon.net/tag/02VF}{Tag 02VF}]{kerodon} the colimit is isomorphic to the space defined by the category $\cat{C}=\int_{(\Fins)\op} F$ whose:
	\begin{itemize}
		\item objects are pairs $(\CI,f:\CI\to A)$ for $\CI\in \Fins$,
		\item morphisms $(\CI,f)\to (\CJ,g)$ are surjections $\sigma:\CJ\surjto \CI$ such that
			\[
			g = f\circ\sigma.
			\] 
	\end{itemize}
	Note that that $\cat{C}$ decomposes as a coproduct:
	\[
		\cat{C} \isoto \coprod_{I\subset \Ran A}\cat{C}_I
	\] 
	where $(\CI,f)$ belong to $\cat{C}_I$ if and only if $\img(f) = I$. But now each $\cat{C}_I$ has an initial object $(I,\id_I)$, so $\cat{C}_I$ is contractible. Hence
	\[
		\cat{C} \simeq \coprod_{I\in \Ran A} \{*\} \isoto \Ran A.
	\] 
\end{proof}

\subsection{\texorpdfstring{$\ell$-adic Sheaves on Prestacks}{l-adic Sheaves on Prestacks}}
Recall the ``universal property of presheaves of spaces'' \cite[\href{https://kerodon.net/tag/03WH}{Tag 03WH}]{kerodon}: left Kan extension defines an equivalence
\[
	\Fun(\Sch,\cat{C}) \simto \Fun^{L}(\PreStk,\cat{C})
\] 
where $\Fun^{L}$ are colimit preserving functors. Hence taking $\cat{C}=\iCat\op$, reversing arrows and right Kan extending $\Shv$ gives us a limit preserving functor
\[
\Shv^{*}:\PreStk\op \to \iCat.
\] 
So we have defined the category of $\ell$-adic sheaves on any prestack! Now we just need to apply this to the Ran space to get
\[
\Shv^{*}(\Ran X) \isoto \lim_{\CI\in \Fins} \Shv(X^{\CI}).
\] 
Hold on, we wanted this to be a colimit\dots

But there is a situation when we can identify a limit in $\iCat$ with a colimit:
\begin{proposition}\label{prop:lim-colim}
	Let $\Phi:\bb{I}\to \PrL$ be a functor and let $\Psi:\bb{I}\op\to \iCat$ be the functor defined as follows:
	\begin{enumerate}[(i)]
		\item On objects $\Phi$ and $\Psi$ agree.
		\item On morphisms $\Psi(\alpha\op)$ is the right adjoint to $\Phi(\alpha)$.
	\end{enumerate}
	Then
	\[
	\colim \Phi \isoto \lim \Psi
	\] 
	in $\PrL$.
\end{proposition}
\begin{proof}[Informal Argument]
	Let $L_i:\Phi(i)\to \cat{C}$ be a compatible family of morphisms §in $\PrL$. Taking right adjoints we get a compatible family of morphisms $R_i:\cat{C}\to \Psi(i)$ in $\PrR$. Now by the universal property of the limit, there is a unique morphism $R:\cat{C}\to \lim \Psi$ in $\PrR$ compatible with $R_i$. Taking the left adjoint we get a unique morphism $L:\lim \Psi\to \cat{C}$ in $\PrL$ compatible with $L_i$. So $\lim \Psi$ satisfies the universal property of the colimit. Finally, since the inclusion $\PrR\injto\iCat$ \cite[Proposition~5.5.3.18]{htt} preserves limits we are done.
\end{proof}
\begin{proof}[Proof]
	The functor $\Psi$ is defined by using \cite[Proposition~5.5.3.4]{htt}. From the same proposition it follows that the limit of $\Psi$ in $\PrR$ is the colimit of $\Phi$ in $\PrL$. The final step of the argument is as in the informal proof.
\end{proof}

So as long as for a morphism $f$ of schemes $\Shv(f)$ is a right adjoint we will get what we want! But for a closed embedding $f$, the pullback $f^{*}$ is not a right adjoint. However, $f^{!}$ is, so we can get the right notion of sheaves on $\Ran X$ by modifying $\Shv$ to send $f\mapsto f^{!}$. Hence we define:
\begin{definition}
	Consider the functor
	\begin{align*}
		\Shv: \Sch\op &\longrightarrow \iCat \\
		S &\longmapsto \Shv(S) \\
		(f:S\from S') &\longmapsto (f^{!}:\Shv(S)\to \Shv(S')).
	\end{align*}
	Let $\Shv^{!}:\PreStk\op\to \iCat$ be its right Kan extension. If $\CY$ is a prestack, then $\Shv^{!}(\CY)$ is the \defn{$\infty$-category of $\ell$-adic sheaves} on $\CY$.
\end{definition}

For a general prestack $\CY$ we have
\[
\Shv^{!}(\CY) \isoto \lim_{S\in \Sch_{/\CY}} \Shv(S).
\] 
Since the inclusion $\PrL\injto \iCat$ preserves limits, we can compute this limit inside $\PrL$. Even better, we can compute it in the full subcategory $\iCatSt$ of $\PrL$ spanned by stable, presentable infinity categories. In particular, $\Shv^{!}$ is endowed with a symmetric monoidal structure, which defines symmetric monoidal structures on $\Shv^{!}(\CY)$ for $\CY\in \PreStk$.

We call prestacks for which \Cref{prop:lim-colim} applies pseudo-schemes:
\begin{definition}
	Let $\CY$ be a prestack. We say $\CY$ is a \defn{pseudo-scheme} if and only if it's expressible as a colimit of schemes with proper transition morphisms.
\end{definition}
\begin{corollary}
	Let
	\[
		\CY \isoto \colim_{i\in \cat{I}} S_i
	\] 
	be a pseudo-scheme. Then
	\[
		\Shv^{!}(\CY) \isoto \colim_{i\in \cat{I}}\Shv(S_i)
	\] 
	with transition maps given by $!$-pushforward.
\end{corollary}
\begin{remark}
	For prestacks, such as ind-schemes or classifying stacks, where we have a good idea of what $\ell$-adic sheaves are supposed be this construction recovers the right thing.
\end{remark}

\subsection{\texorpdfstring{$\ell$-adic Sheaves on the Ran Space}{l-adic Sheaves on the Ran Space}}
Let us unwind a sheaf $\CF\in \Shv^{!}(\Ran X)$ looks like.

The limit description tells us that $\CF$ is the data of:
\begin{itemize}
	\item For every scheme $S\in \Sch$ and non-empty, finite subset $I\subset X(S)$ an $\ell$-adic sheaf $\CF_{S,I}$ on $S$.
	\item For every morphism of schemes $f:S'\to S$ and non-empty, finite subset $I\subset X(S)$ an isomorphism
		\[
		\CF_{S',f(I)} \simto f^{!}\CF_{S,I}.
		\] 
\end{itemize}
These isomorphism then need to satisfy higher coherences.

The colimit description tells us that $\CF$ is the data of:
\begin{itemize}
	\item For every non-empty, finite set $\CI$ an $\ell$-adic sheaf $\CF_\CI$ on $X^{\CI}$.
	\item For every surjective map $\CJ\surjto \CI$ of sets an isomorphism
		\[
		f_!\CF_{\CI} \simto \CF_{\CJ},
		\] 
		where $f:X^{\CI}\to X^{\CJ}$ is the morphism corresponding to $\CJ\surjto \CI$.
\end{itemize}
These isomorphisms need to satisfy higher coherences.

\subsection{Homology}
Derived pushforwards are perhaps the most important tool for computations in cohomology, so if we want to get any decent understanding of $\Shv^{!}(\Ran X)$ it would be good to have access to them. On its face it's not at all obvious that for a morphism $f:\CY_1\to \CY_2$ either $f_*$ or $f_!$ should exist.

First, note that we are given $f^{!}$ for free by the definition of $\Shv^{!}$.

\begin{proposition}
	The right adjoint $f_!$ to $f^{!}$ exists.
\end{proposition}
\begin{proof}
	By the Adjoint Functor Theorem \cite[Corollary~5.5.2.9(ii)]{htt} it suffices to show that $f^{!}$ preserves limits. This is the content of \Cref{lmm:pres-lim}.
\end{proof}

\begin{lemma}\label{lmm:pres-lim}
	$f^{!}$ preserves limits.
\end{lemma}
\begin{proof}
	Note that for $(S,y)\in \Sch_{/\CY}$ and a family $\CF^{i}\in \Shv^{!}(\CY)$ indexed by $\cat{I}$ we have
	\[
		\left( \lim_{i\in \cat{I}}\CF^{i} \right)_{S,y} \isoto \lim_{i\in \cat{I}} \CF^{i}_{S,y}.
	\] 
	But exceptional pullback for schemes preserves limits, so we are done.
\end{proof}

This allows us to define the \emph{homology} of a prestack $\CY$ as follows. Taking $\CY_1=\CY$ and $\CY_2=\Spec k$ we get a functor
\begin{align*}
	f_!: \Shv^{!}(\CY) &\longrightarrow \lMod \\
	\CF &\longmapsto C_c^{*}(\CY,\CF).
\end{align*}
\begin{definition}
	Let $\CY$ be a prestack, $p_\CY:\CY\to \Spec k$ be the structure morphism and $p_\CY^{!}:\lMod\to \Shv^{!}(\CY)$ be the functor defined by $p_\CY$. The \defn{dualizing sheaf} on $\CY$ is
	\[
	\omega_\CY \defeq p_\CY^{!}\Lambda.
	\] 
\end{definition}
\begin{definition}
	The \defn{homology} of a prestack $\CY$ is
	\[
	C_*(\CY) \defeq C^{*}_c(\CY,\omega_\CY).
	\] 
	The \defn{reduced homology} of $\CY$ is
	\[
	C_*\redu(\CY) \defeq \Fib(C_*(\CY)\to C_*(\Spec k)) = \Fib \left( p_{\CY,!}p_\CY^{!}\Lambda \to \Lambda \right) .	
	\] 
\end{definition}

We actually know what the homology of the Ran space is!
\begin{theorem}[Beilinson--Drinfeld]
	Suppose $X$ is separeted and connected. Then the structure morphism $\Ran X\to \Spec k$ induces a quasi-isomorphism
	\[
	C_*(\Ran X) \simto C_*(\Spec k).
	\] 
\end{theorem}
\begin{proof}
	See \cite[Theorem~2.4.5]{weil}.
\end{proof}
\begin{corollary}
	We have canonical isomorphisms
	\[
	H_q(\Ran X) \isoto \begin{cases}
		\Lambda & \text{if }q=0 \\
		0 & \text{otherwise}.
	\end{cases}
	\] 
\end{corollary}

\subsection{Computing the Pushforward I}
The next natural step in understanding the Ran space is to compute
\begin{equation}\label{eq:ins2}
\ins_\CJ^{!}\ins_{\CI,!} \CF
\end{equation} 
for $\CI,\CJ\in \Fins$ and $\CF\in \Shv(X^{\CI})$. For a general morphism $f$ of prestacks $f_!$ is intractable, but fortunately for prestacks similar to the Ran space we can do better.

But first, we show that in the case of the Ran space understanding \eqref{eq:ins2} in a real sense describes all of the sheaves on it.

\begin{lemma}
	Let
	\[
		\CY = \colim_{i\in \cat{I}}Z_i
	\] 
	be a pseudo-scheme. Then $\ins_i^{!}$ preserve colimits.
\end{lemma}
\begin{proof}
	This is because colimits in $\Shv^{!}(\CY)$ are computed pointwise.
\end{proof}
\begin{corollary}
	Let
	\[
		\CY = \colim_{i\in \cat{I}}Z_i
	\] 
	be a pseudo-scheme and $\CF\in \Shv^{!}(Y)$. Then the natural map
	\[
		\colim_{i\in \cat{I}} \ins_{i,!}\ins_i^{!}\CF \to \CF
	\] 
	is an isomorphism.
\end{corollary}

Now we suss out the properties of $\Ran X$ that allow us to compute $\ins_{\CI,!}$.

\begin{definition}
	A morphism $f:\CY\to S$ of prestacks with $S\in \Sch$ is \defn{pseudo-proper} if $\CY$ can be expressed as a colimit of proper schemes over $S$.
\end{definition}
\begin{proposition}
	Let
	\[
	\begin{tikzcd}
		\CY' \ar[r,"g_\CY"] \ar[d,"f'"] & \CY \ar[d,"f"] \\
		S' \ar[r,"g_S"] & S
	\end{tikzcd}
	\] 
	be a Cartesian diagram in $\PreStk$ with $S,S'\in \Sch$ and $f$ pseudo-proper. Then the natural morphism
	\[
	f'_!\circ g_\CY^{!} \to g_S^{!} \circ f_!
	\] 
	is an isomorphism.
\end{proposition}
\begin{proof}
	Follows from Proper Base Change for schemes.
\end{proof}
\begin{definition}
	A morphism $f:\CY_1\to \CY_2$ of prestacks is \defn{pseudo-proper} if for every morphism $S\to \CY_2$ with $S\in \Sch$ the morphism $f_S:\CY_1\times_{\CY_2}S\to S$ is pseudo-proper.
\end{definition}
\begin{corollary}\label{cor:bc}
	Let
	\[
	\begin{tikzcd}
		\CY_1' \ar[r,"g_1"] \ar[d,"f'"] & \CY_1 \ar[d,"f"] \\
		\CY_2' \ar[r,"g_2"] & \CY_2
	\end{tikzcd}
	\] 
	be a Cartesian diagram in $\PreStk$ with $f$ pseudo-proper. Then the natural morphism
	\[
	f'_!\circ g_\CY^{!} \to g_S^{!} \circ f_!
	\] 
	is an isomorphism.
\end{corollary}

So we would like to verify that $\ins_\CI$ are pseudo-proper.
\begin{proposition}\label{prop:insp}
	Let
	\[
		\CY = \colim_{i\in \cat{I}}Z_i
	\] 
	be a pseudo-scheme. Then $\ins_i$ are pseudo-proper.
\end{proposition}
\begin{proof}
	A bit of abstract nonsense: see \cite[Proposition~7.4.2]{abott}.
\end{proof}

Let $\CI,\CJ\in \Fins$ and define the prestack $\wt{X}^{\CI,\CJ}$ by the following Cartesian diagram
\[
\begin{tikzcd}
	\wt{X}^{\CI,\CJ} \ar[r,"\wt{q}_{\CI}"] \ar[d,"\wt{q}_{\CJ}"] & X^{\CI} \ar[d,"\ins_{\CI}"] \\
	X^{\CJ} \ar[r,"\ins_{\CJ}"] & \Ran X
\end{tikzcd}
\] 
Now by \Cref{prop:insp} and \Cref{cor:bc}
\[
	\ins_\CJ^{!}\circ\ins_{\CI,!} \isoto \wt{q}_{\CJ,!}\circ \wt{q}_{\CI}^{!}.
\] 
Unfortunately, $\wt{X}^{\CI,\CJ}$ is still not a scheme. But what if we could find a scheme $X^{\CI,\CJ}$ with an equivalent category of sheaves?

\subsection{Computing the Pushforward II}
\paragraph{\textbf{Disclaimer:} I couldn't make sense of the paper here, so this section is my feeble attempt at finding my own argument.}

We can describe $\wt{X}^{\CI,\CJ}$ as a colimit
\[
	\wt{X}^{\CI,\CJ} \isoto \colim_{\CI\surjto\CK\surjfrom\CJ} X^{\CK}
\] 
and the morphism to $X^{\CI}\times X^{\CJ}$ is the obvious one. The morphism to $X^{\CI}\times X^{\CJ}$ satisfies some important properties.
\begin{definition}
	Let $S\in \Sch$ and
	\[
		\CY = \colim_{i\in \cat{I}}Z_i
	\] 
	for $Z_i\in \Sch_{/S}$. We say that $\CY$ is:
	\begin{enumerate}[(i)]
		\item a \defn{finitary} pseudo-scheme over $S$ if $\cat{I}$ is finite;
		\item \defn{pseudo-proper} over $S$ if $Z_i$ are prper over $S$;
		\item \defn{finitary pseudo-proper} over $S$ if it is a finitary pseudo-scheme that's pseudo-proper over $S$.
	\end{enumerate}
\end{definition}
Since $X^{\CK}\injto X^{\CI}\times X^{\CJ}$ are closed embeddings, $\wt{X}^{\CI,\CJ}$ is finitary pseudo-proper over $X^{\CI}\times X^{\CJ}$.

Now define
\[
	\wt{X}^{\CI,\CJ}|\red \defeq \colim_{\CI\surjto \CK \surjfrom \CJ}X^{\CK}|\red.
\] 
Since the etale topology doesn't see nilpotents, the inclusion
\[
	\wt{X}^{\CI,\CJ}|\red \injto \wt{X}^{\CI,\CJ}
\] 
induces an equivalence
\[
	\Shv^{!}(\wt{X}^{\CI,\CJ}) \simto \Shv^{!}(\wt{X}^{\CI,\CJ}).
\] 

Now let $X^{\CI,\CJ}$ be the reduced subscheme defined by the union of images of $X^{\CK}$, then the morphism
\[
	\wt{X}^{\CI,\CJ}|\red \to X^{\CI}\times X^{\CJ}
\] 
factors through $X^{\CI,\CJ}$. In fact,
\[
	g:\wt{X}^{\CI,\CJ}|\red \to X^{\CI,\CJ}
\] 
is finitary pseudo-proper, surjective on $k$-points and injective on $S$-points for $S\in \Sch$. Now we use
\begin{lemma}
	Let $f:\CY_1\to Y_2$ be a pseudo-proper morphism of prestacks with $Y_2\in \Sch$.
	\begin{enumerate}[(a)]
		\item If f is \defn{injective}, i.e. $f(S)$ is a monomorphism for all $S\in \Sch$, then $f_!$ is fully faithful.
		\item If $f$ is finitary pseudo-proper and surjective on $k$-points, then $f^{!}$ is conservative.
	\end{enumerate}
\end{lemma}
\begin{proof}
	\cite[Lemma~7.4.11]{abott}.
\end{proof}
Hence we see that $(g_!,g^!)$ define an equivalence of categories. Hence finally, writing
\[
\begin{tikzcd}
	X^{\CI} & X^{\CI,\CJ} \ar[l,"q_{\CI}"] \ar[r,"q_{\CJ}"] & X^{\CJ}
\end{tikzcd}
\] 
we find that:
\begin{corollary}
	We have canonical isomorphisms
	\[
		\ins_\CJ^{!}\circ \ins_{\CI,!} \isoto \wt{q}_{\CJ,!}\circ \wt{q}_{\CI}^{!} \isoto q_{\CJ,!}\circ q_{\CI}^{!}.
	\] 
\end{corollary}

\appendix
\printbibliography

\end{document}
